\documentclass[a4paper,10pt,twoside]{report}
\usepackage[a4paper]{geometry}
\usepackage{listings}
\usepackage[latin1]{inputenc} % ISO 8859-1
\usepackage[ngerman]{babel}
\usepackage{color}
\usepackage[pdftex]{graphicx}
\usepackage{fancyhdr}
\usepackage{amsmath,amssymb,amsfonts,textcomp}
\usepackage{hyperref}
\usepackage[compact]{titlesec} % Kapitel�berschriften
\usepackage[T1]{fontenc}
\usepackage{xspace}



% Dateipfade
\graphicspath{{pictures/}}

% Eigene Farben
\definecolor{lightgray}{rgb}{0.95,0.95,0.95}
\definecolor{codegray}{rgb}{0.8,0.8,0.8}
\definecolor{darkgray}{rgb}{0.18,0.18,0.18}
\definecolor{mygreen}{rgb}{0.133,0.545,0.133}
\definecolor{mypurple}{rgb}{0.627,0.126,0.941}

% PDF Informationen
\title{Beispieldokumentation zum Projektpraktikum}
\author{Hier bitte Autor eintragen}

% Seitengestaltung
\setlength\textheight{25.0cm}
\setlength\topmargin{-1.3cm}
\setlength\textwidth{1.1\textwidth}
% ohne Heftkante
\setlength\oddsidemargin{0pt}
\setlength\evensidemargin{0pt}
% mit Heftkante
%\setlength\oddsidemargin{+10pt}
%\setlength\evensidemargin{-10pt}
\setlength{\headheight}{13pt}

% Kopf- und Fußzeilen
\pagestyle{fancy}
\fancypagestyle{plain}{\thispagestyle{fancy}}

\fancyhead{}
%\fancyheadoffset[LE,RO]{0pt}
\fancyhead[LO,RE]{\rightmark} % aktuelle Section
\fancyhead[RO,LE]{\includegraphics[viewport=0 7.5 168 24,scale=0.5]{header}} % Logo here

\fancyfoot{}
%\fancyfootoffset[LE,RO]{0pt}
\fancyfoot[LO,RE]{\leftmark} % aktuelle Chapter
\fancyfoot[LE]{\thepage\hspace{3em}Dokumentation zum IT-Projektpraktikum} %Seitenzahl links
\fancyfoot[RO]{Dokumentation zum IT-Projektpraktikum\hspace{3em}\thepage} %Seitenzahl recht

% Titel Linie
\newcommand{\bigrule}{\titlerule[0.5mm]}

% Titel Format
\titleformat{\chapter}[display]
{\bfseries\Huge}
{%
 \vskip-2em
 \titlerule
 \filright
 \Large\chaptertitlename\
 \Large\thechapter
}
{0mm}
{\filright}
[\bigrule]

% Definitionen für Listing
\lstset{ %
language=C++,                % choose the language of the code
numbers=left,                   % where to put the line-numbers
%numberstyle=\footnotesize,      % the size of the fonts that are used for the line-numbers
stepnumber=1,                   % the step between two line-numbers. If it's 1 each line will be numbered
numbersep=5pt,                  % how far the line-numbers are from the code
backgroundcolor=\color{white},  % choose the background color. You must add \usepackage{color}
showspaces=false,               % show spaces adding particular underscores
showstringspaces=false,         % underline spaces within strings
%showtabs=false,                 % show tabs within strings adding particular underscores
frame=single,                   % adds a frame around the code
tabsize=2,                      % sets default tabsize to 2 spaces
captionpos=t,                   % sets the caption-position to bottom
breaklines=true,                % sets automatic line breaking
xleftmargin=3.5pt,
xrightmargin=3.5pt,
breakatwhitespace=false,        % sets if automatic breaks should only happen at whitespace
extendedchars=true,captionpos=t , prebreak=\mbox{$\hookleftarrow$}, belowcaptionskip=0.5em,postbreak={},
basicstyle=\ttfamily,       % the size of the fonts that are used for the code
keywordstyle=\color{blue}\bfseries\sffamily,
commentstyle=\color{mygreen}\slshape,
stringstyle=\color{mypurple},
}

\setlength{\parindent}{0pt} % Texteinzug bei neuem Absatz

% Zähler für Ausgabe
\newcounter{ausgabe}[chapter]
\renewcommand\theausgabe{\thechapter.\arabic{ausgabe}}

% Programmausgabe, grau hinterlegen
\newcommand{\ausg}[2]{
	\stepcounter{ausgabe}

	\vspace{2mm}
	\begin{minipage}{\textwidth}
	  \colorbox{codegray}{
			\begin{minipage}{0.97\textwidth}
				\textbf{\texttt{#2}}
			\end{minipage}
		}
		\begin{center}
			Ausgabe \theausgabe: #1
		\end{center}
	\end{minipage}
}

% Schreibt 'Code::Blocks'
\newcommand{\cb}{\mbox{Code::Blocks}\xspace}

% Schreibt 'Visual Studio'
\newcommand{\vs}{\mbox{Visual Studio}\xspace}

% TODO - Command
\newcommand{\TODO}[1]{{\color{red}TODO!! #1}}

% code - Command
\newcommand{\code}[1]{
\newpage{#1}
\lstinputlisting[caption={#1}]{#1}
}

% Leerseite ohne Seitennummer, nächste Seite rechts
\newcommand{\blankpage}{
 \clearpage{\pagestyle{empty}\cleardoublepage}
}



\begin{document}
%----------------------------
\begin{titlepage}
	\begin{flushleft}
		\begin{picture}(0,0)
			\put(0,-2){\includegraphics[height=4em]{KIT_Logo.pdf}}
		\end{picture}
			\put(140,14){\Large Karlsruher Institut f�r Technologie}
			\put(120,0){\large Institut f�r Technik der Informationsverarbeitung}
		\begin{picture}(0,0)
			\put(410,-2){\includegraphics[height=4em]{itiv}}
		\end{picture}
	\end{flushleft}

	\begin{center}
		\vspace{7cm}
		\addvspace{1.0cm}
		\Huge {Dokumentation zum Projektpraktikum Informationstechnik}\\
	\end{center}
	\vspace{4.0cm}
	\Large

	\begin{flushleft}
	\begin{tabular}{ll}
		Gruppe:& 064\\
		Gruppenmitglieder:& Benedikt Braunger, Cornelius Richt,\\ & Fabian Schackmar, Lukas Mohrbacher\\
		Tutor:& Florian Brauchle\\
		Abgabetermin:& 18.01.2013\\
		Semester:& WS2011/2012
	\end{tabular}

	\vfill

	KIT - Universit�t des Landes Baden W�rttemberg und nationales Forschungszentrum in der Helmholtz-Gemeinschaft
	\end{flushleft}
\end{titlepage}
%----------------------------
\tableofcontents
\lstlistoflistings
\appendix
\chapter{Quelltext}
\code{main.cpp}
\code{Menue.h}
\code{Menue.cpp}
\code{Faktoren.h}
\code{Faktoren.cpp}
\code{Bibliothek.h}
\code{Bibliothek.cpp}
\code{GatterTyp.h}
\code{GatterTyp.cpp}
\code{Flipflop.h}
\code{Flipflop.cpp}
\code{SignalListeErzeuger.h}
\code{SignalListeErzeuger.cpp}
\code{signals.h}
\code{signals.cpp}
\code{SchaltwerkElement.h}
\code{SchaltwerkElement.cpp}
\code{ListenElement.h}
\code{ListenElement.cpp}
\code{GraphErzeuger.h}
\code{GraphErzeuger.cpp}
\code{LaufzeitAnalysator.h}
\code{LaufzeitAnalysator.cpp}
\code{cross-compatibility.h}
\code{cross-compatibility.cpp}

\end{document}
